% ============================================================
% CAPITOLUL 7 — COMPARAȚIE, ANALIZĂ NUMERICĂ ȘI CALIBRARE
% ============================================================

\chapter{Comparația Modelelor, Analiză Numerică și Calibrare}
\label{ch:comparatie_calibrare}

\section{Rolul comparației între modele}
\label{sec:rol_comparatie_modele}

Modelele compartimentale discutate în lucrare nu trebuie interpretate ca variante
interschimbabile ale aceleiași ecuații. Fiecare model introduce o anumită
ipoteză epidemiologică și elimină alte detalii pentru a păstra sistemul
analizabil. Din acest motiv, comparația trebuie făcută pe două niveluri:
compararea structurală a ipotezelor și compararea numerică a traiectoriilor
obținute pe același set de date.

Comparația structurală răspunde la întrebarea „ce fenomen poate reprezenta
modelul?”. Comparația numerică răspunde la întrebarea „cât de bine reproduce
modelul datele observate în condiții experimentale identice?”. Ambele sunt
necesare. Un model simplu poate avea o eroare mică pe o fereastră scurtă, dar
poate ascunde mecanisme importante; un model complex poate reprezenta mai fidel
fenomenul, dar poate fi dificil de calibrat din date puține.

\section{Comparație structurală}
\label{sec:comparatie_structurala}

Tabelul~\ref{tab:comparatie_modele} sintetizează diferențele principale dintre
modelele suportate de aplicație. Criteriile alese sunt cele relevante pentru
implementarea curentă: compartimentele, parametrii necesari, ipoteza centrală și
tipul de date observabile cu care modelul poate fi comparat direct.

\begin{table*}[t]
\centering
\caption{Comparație structurală între modelele utilizate în lucrare.}
\label{tab:comparatie_modele}
\begin{tabularx}{0.96\textwidth}{l l l X X}
\toprule
\textbf{Model} & \textbf{Compartimente} & \textbf{Parametri} & \textbf{Ipoteză principală} & \textbf{Observabil comparabil} \\
\midrule
SIS & $S,I$ & $\beta,\gamma$ & infecția nu produce imunitate durabilă & cazuri active aproximabile prin $I$ \\
SIR & $S,I,R$ & $\beta,\gamma$ & recuperarea produce imunitate permanentă & active prin $I$, cumulate prin $I+R$ \\
SEIR & $S,E,I,R$ & $\beta,\sigma,\gamma$ & există o fază latentă neinfecțioasă sau slab observabilă & active prin $I$, latența rămâne neobservată \\
SIRS & $S,I,R$ & $\beta,\gamma,\xi$ & imunitatea se pierde în timp & active prin $I$, revenire din $R$ în $S$ \\
SIDARTHE & $S,I,D,A,R,T,H,E$ & 16 rate & diagnosticul, simptomatologia și severitatea clinică modifică dinamica observată & active diagnosticate prin $D+R+T$, decese prin $E$ \\
\bottomrule
\end{tabularx}
\end{table*}

SIS este potrivit pentru boli unde reinfectarea rapidă sau lipsa imunității este
o ipoteză acceptabilă. SIR este util ca reper matematic, deoarece păstrează
dinamica minimală a unei epidemii cu imunitate. SEIR introduce latența și devine
mai potrivit atunci când există o perioadă de incubație importantă. SIRS relaxează
ipoteza imunității permanente. SIDARTHE mută accentul de la agregarea biologică
la observabilitatea clinică: infecțiile nedepistate, cazurile diagnosticate,
formele simptomatice, formele severe, vindecările și decesele sunt tratate
separat.

\section{Operatorul de observație}
\label{sec:operator_observatie}

Datele reale nu conțin de obicei toate compartimentele modelului. Pentru a compara
un model cu o serie observată este necesar un operator de observație
$g_m(\mathbf{x}(t))$, dependent de modelul $m$, care transformă starea internă a
modelului într-o mărime comparabilă cu datele.

Pentru o serie de cazuri active, operatorul poate fi definit astfel:
\begin{equation}
g_{\text{SIR}}(\mathbf{x}) = I, \qquad
g_{\text{SEIR}}(\mathbf{x}) = I, \qquad
g_{\text{SIRS}}(\mathbf{x}) = I.
\end{equation}
Pentru SIDARTHE, dacă seria observată reprezintă cazuri diagnosticate active, o
alegere mai realistă este
\begin{equation}
g_{\text{SIDARTHE}}(\mathbf{x}) = D + R + T,
\end{equation}
deoarece compartimentul $I$ reprezintă infectați nedepistați, iar $A$ reprezintă
persoane simptomatice nedepistate. Dacă setul de date conține doar o aproximare
generală a cazurilor active, aplicația poate folosi o proiecție mai simplă, dar
interpretarea trebuie notată explicit.

Pentru decese cumulative, operatorul SIDARTHE este direct:
\begin{equation}
g_{\text{decese}}(\mathbf{x}) = E.
\end{equation}
Această separare este un avantaj al modelului SIDARTHE față de modelele SIR,
SEIR sau SIRS, în care decesele sunt incluse implicit în compartimentul
\emph{removed}.

\section{Metrici de eroare numerică}
\label{sec:metrici_eroare_numerica}

Fie $y_t$ seria observată la ziua $t$, iar $\hat{y}_{m,t}$ seria produsă de
modelul $m$ după aplicarea operatorului de observație. Pentru o fereastră de
evaluare cu $T$ puncte, eroarea medie absolută este
\begin{equation}
\operatorname{MAE}_m =
\frac{1}{T}\sum_{t=1}^{T}\left|y_t-\hat{y}_{m,t}\right|.
\end{equation}
Eroarea medie pătratică radicală este
\begin{equation}
\operatorname{RMSE}_m =
\sqrt{\frac{1}{T}\sum_{t=1}^{T}\left(y_t-\hat{y}_{m,t}\right)^2}.
\end{equation}
Pentru a compara regiuni cu populații diferite, este utilă normalizarea la
populația totală $N$:
\begin{equation}
\operatorname{nRMSE}_m = \frac{\operatorname{RMSE}_m}{N}.
\end{equation}

În date epidemiologice, MAPE poate fi instabilă atunci când $y_t$ este aproape
de zero. Din acest motiv, lucrarea preferă MAE, RMSE și nRMSE. Dacă se compară
mai multe serii observabile, de exemplu cazuri active și decese, se folosește o
funcție de pierdere ponderată:
\begin{equation}
\mathcal{L}_m(\theta) =
\sum_{k \in \mathcal{K}} w_k
\frac{1}{T}\sum_{t=1}^{T}
\left(y_{k,t}-g_{m,k}(\mathbf{x}_{\theta}(t))\right)^2,
\end{equation}
unde $k$ indexează tipul de observație, iar $w_k$ controlează importanța fiecărei
serii. Ponderile sunt necesare deoarece seriile pot avea scări diferite.

\testtodo{Formulele de mai sus definesc modul de calcul al erorilor, dar nu
reprezintă rezultate în sine. Pentru a obține rezultate numerice trebuie rulate
toate modelele pe același set de date, trebuie aplicat operatorul de observație ales
și trebuie raportate valorile MAE, RMSE și nRMSE pentru fiecare combinație
model--limbaj--fereastră temporală.}

\section{Protocol de analiză comparativă}
\label{sec:protocol_comparativ}

Pentru ca analiza comparativă să fie reproductibilă, toate modelele trebuie
rulate în condiții identice. Protocolul propus are următorii pași:

\begin{enumerate}
    \item se alege același set de date, aceeași regiune și aceeași fereastră
    temporală pentru toate modelele;
    \item se aplică aceeași regulă de inițializare a compartimentelor, conform
    Tabelului~\ref{tab:initializare_modele};
    \item se calibrează fiecare model pe aceeași fereastră de antrenare;
    \item se evaluează eroarea pe o fereastră separată de test;
    \item se raportează MAE, RMSE, nRMSE și timpul de execuție;
    \item se păstrează separat interpretarea epidemiologică și costul software.
\end{enumerate}

Separarea între antrenare și test este importantă. Un model cu mulți parametri,
cum este SIDARTHE, poate potrivi mai bine fereastra de antrenare, dar aceasta nu
garantează o predicție mai bună pe zile viitoare. Prin urmare, analiza
comparativă trebuie să raporteze atât eroarea de calibrare, cât și eroarea de
generalizare.

\begin{table*}[t]
\centering
\caption{Locurile din analiza comparativă care necesită testare pentru obținerea rezultatelor.}
\label{tab:teste_necesare_comparatie}
\begin{tabularx}{0.96\textwidth}{lX X}
\toprule
\textbf{Secțiune} & \textbf{Test necesar} & \textbf{Rezultat de raportat} \\
\midrule
Operatorul de observație & alegerea proiecției potrivite pentru fiecare model și set de date & seria comparabilă cu datele reale: active, diagnosticate sau decese \\
Metrici de eroare & rularea modelelor pe aceeași fereastră temporală & MAE, RMSE, nRMSE pe fiecare model și limbaj \\
Protocol comparativ & împărțirea cronologică train/test și rulări identice pentru toate modelele & tabel agregat cu eroare pe train și test \\
Calibrare parametri & optimizarea parametrilor în limite admisibile & parametrii calibrați, eroarea finală și numărul de evaluări ale solverului \\
Identificabilitate & repetarea calibrării din puncte inițiale diferite & stabilitatea parametrilor și variația erorii între rulări \\
\bottomrule
\end{tabularx}
\end{table*}

\section{Calibrarea parametrilor}
\label{sec:calibrare_parametri}

Calibrarea parametrilor transformă modelul dintr-un instrument de simulare
exploratorie într-un instrument ajustat la date observate. Pentru un model
$m$ cu vector de parametri $\theta$, calibrarea poate fi formulată ca problemă
de optimizare:
\begin{equation}
\hat{\theta}_m =
\arg\min_{\theta \in \Omega_m} \mathcal{L}_m(\theta),
\end{equation}
unde $\Omega_m$ este domeniul admisibil al parametrilor. Constrângerile sunt
esențiale: ratele de tranziție trebuie să fie nenegative, iar intervalele lor
trebuie alese în funcție de scară temporală și de cunoștințe epidemiologice.

Pentru modelele SIR/SIS/SEIR/SIRS, calibrarea poate fi făcută cu metode clasice
de optimizare neliniară, de exemplu Nelder--Mead \cite{nelder1965simplex},
L-BFGS-B sau căutare pe grilă pentru intervale mici. Pentru SIDARTHE, spațiul de
parametri este mai mare, iar optimizarea devine mai fragilă. În acest caz sunt
utile:
\begin{itemize}
    \item fixarea unor parametri din literatură sau din ipoteze clinice;
    \item calibrarea pe grupuri de parametri, nu a tuturor ratelor simultan;
    \item regularizarea față de valori de referință;
    \item validarea pe o fereastră temporală nefolosită la optimizare.
\end{itemize}

\subsection{Aplicare pe Romania, valul 1 din 2020}
\label{subsec:calibrare_romania}

Procedura este aplicată concret modelului SIR pe seria de cazuri active din
Romania (sursa OWID, vezi capitolul~\ref{ch:implementare}), pe valul epidemic
de primăvară-vară 2020. Implementarea este în
\texttt{scripts/calibrate.py} și folosește
\texttt{scipy.optimize.minimize} cu metoda Nelder--Mead, restartată din opt
puncte inițiale aleatoare pentru a reduce sensibilitatea la condiția
inițială. Fereastra este împărțită cronologic: $15$~martie -- $15$~august 2020
($154$ zile) pentru antrenare, $16$~august -- $15$~septembrie 2020 ($31$ zile)
pentru test. Funcția de pierdere este eroarea pătratică medie pe seria de
cazuri active, iar operatorul de observație este $g_{\text{SIR}}(\mathbf{x})=I$.

Calibrarea este executată în două variante:
\begin{itemize}
    \item \emph{liber}: optimizare simultană a $(\beta, \gamma)$ în
    $\Omega = (0, 5] \times (0, 5]$;
    \item \emph{constrans}: $\gamma$ fixat la $1/14$ zi$^{-1}$, valoarea de
    referință din literatura clinică pentru COVID-19 (durata medie a perioadei
    infecțioase $\sim 14$ zile), iar optimizarea se face doar pe $\beta$.
\end{itemize}
Comparația este prezentată în Tabelul~\ref{tab:calibration_results} și
Figura~\ref{fig:calibration_sir_fit}.

\begin{table*}[!htbp]
\centering
\caption[Calibrare SIR pe România]{Calibrare SIR pe România, valul 1 ($15$~martie -- $15$~august 2020 antrenare, $16$~august -- $15$~septembrie 2020 test). Două scenarii: \emph{liber} (optimizează ambele rate, $(\beta,\gamma)$), \emph{constrâns} (folosește $\gamma=1/14$ din durata clinică medie a infecției COVID-19, optimizează doar $\beta$). Comparația ilustrează tensiunea dintre minimizarea pierderii și interpretabilitatea epidemiologică a parametrilor.}
\label{tab:calibration_results}
\begin{tabularx}{\textwidth}{l r r X}
\toprule
\textbf{Marime} & \textbf{Liber} & \textbf{Constrâns} & \textbf{Observatii} \\
\midrule
$\hat{\beta}$ (zi$^{-1}$) & 0.3677 & 0.1039 & rata de transmitere \\
$\hat{\gamma}$ (zi$^{-1}$) & 0.3346 & 0.0714 (\textit{fix}) & $1/\hat{\gamma}$ liber $= 3.0$ zile (non-fizic), constrâns $= 14$ zile \\
$\hat{\mathcal{R}}_0$ & 1.10 & 1.45 & valori de literatură pentru tulpina originala $\sim 2$--$3$ \\
\midrule
MAE antrenare & 1630 & 1612 & cazuri \\
RMSE antrenare & 2199 & 2203 & cazuri \\
MAE test & 13031 & 16727 & extrapolare 1 luna \\
RMSE test & 14889 & 19371 & cazuri \\
\bottomrule
\end{tabularx}
\end{table*}

\begin{figure*}[t]
\centering
\begin{tikzpicture}
\begin{axis}[
    width=0.95\textwidth,
    height=6.5cm,
    xlabel={Zi (de la $15$~martie 2020)},
    ylabel={Cazuri active observate},
    xmin=0, xmax=185,
    ymin=0,
    legend style={at={(0.02, 0.98)}, anchor=north west, font=\footnotesize, fill=white, draw=gray!50},
    grid=both, grid style={line width=.1pt, draw=gray!20},
    minor tick num=1,
    tick label style={font=\footnotesize},
    label style={font=\small},
    legend cell align={left},
]
% bara verticală separare antrenare / test
\draw[dashed, gray!70, thick] (axis cs:154, 0) -- (axis cs:154, 22000)
    node[above, font=\scriptsize, gray!90] {antrenare $|$ test};

% observat (linie continuă)
\addplot[mark=none, thick, black]
    table[col sep=comma, x=day, y=observed]
    {data/calibration_sir_romania.csv};
\addlegendentry{Observat (România OWID)}

% fit liber pe antrenare + forecast pe test
\addplot[mark=none, thick, blue]
    table[col sep=comma, x=day, y=fit_free]
    {data/calibration_sir_romania.csv};
\addlegendentry{Fit liber, $\hat\beta=0{,}368$, $\hat\gamma=0{,}335$}

\addplot[mark=none, thick, dashed, blue!80]
    table[col sep=comma, x=day, y=pred_free]
    {data/calibration_sir_romania.csv};
\addlegendentry{Prognoza liber}

% fit constrâns
\addplot[mark=none, thick, red]
    table[col sep=comma, x=day, y=fit_fixed]
    {data/calibration_sir_romania.csv};
\addlegendentry{Fit constrâns, $\hat\beta=0{,}104$, $\gamma=1/14$}

\addplot[mark=none, thick, dashed, red!80]
    table[col sep=comma, x=day, y=pred_fixed]
    {data/calibration_sir_romania.csv};
\addlegendentry{Prognoza constrâns}

\end{axis}
\end{tikzpicture}
\caption[Calibrare SIR pe valul 1 (2020)]{Calibrare SIR pe România, valul 1 din 2020. Linia neagră arată seria
observată; cele două scenarii produc fituri vizual aproape identice pe
fereastra de antrenare (MAE $\approx 1\,600$ cazuri pentru ambele), dar
prognoza pe luna următoare diverge (MAE $13\,031$ pentru fitul liber vs
$16\,727$ pentru cel constrâns). Parametrii estimați diferă substantial,
ilustrând identifiabilitatea practică slaba: aceleași date sunt explicate
aproape la fel de bine de combinații $(\beta, \gamma)$ foarte diferite.}
\label{fig:calibration_sir_fit}
\end{figure*}


Rezultatele evidențiază două aspecte caracteristice acestei clase de probleme.
În primul rând, eroarea pe antrenare este aproape identică între cele două
scenarii (MAE $\approx 1\,600$ cazuri pentru ambele), deși parametrii estimați
diferă substanțial: $\hat\gamma_{\text{liber}} = 0{,}335$ corespunde unei
durate infecțioase medii de $3$ zile, valoare care nu se susține clinic
pentru COVID-19. Constrângerea $\gamma = 1/14$ produce un $\hat\beta$ de
$\sim 3{,}5$ ori mai mic și un $\hat{\mathcal{R}}_0$ de $1{,}45$, mai apropiat
de intervalul de literatură $2$--$3$ pentru tulpina originală, dar tot sub
această limită din cauza datelor agregate ale unei singure regiuni intr-o
singură fereastră.

În al doilea rând, eroarea pe fereastra de test creste de aproximativ un
ordin de mărime fata de antrenare ($13\,031$ cazuri pentru fitul liber,
$16\,727$ pentru cel constrans). Modelul SIR cu coeficienti constanti nu poate
extrapola al doilea val epidemic, care apare în septembrie-octombrie 2020 cu
$\beta$ efectiv mai mare decât în vară. Aceasta este o limitare structurală a
SIR clasic, nu o problemă de calibrare.

\begin{algorithm}[t]
\caption{Calibrarea unui model compartimental pe date observate}
\label{alg:calibrare_model}
\begin{lstlisting}[style=pseudocode]
Input: model m, dataset Y, parametri initiali theta0,
       limite Omega, fereastra train, operator g_m
Output: parametri calibrati theta_hat, metrici de eroare

theta <- theta0
while criteriul de oprire nu este indeplinit do
    ruleaza solverul ODE cu parametrii theta
    proiecteaza starea modelului prin g_m
    calculeaza functia de pierdere L_m(theta)
    actualizeaza theta prin algoritmul de optimizare
    proiecteaza theta inapoi in Omega daca este necesar
end while
calculeaza MAE, RMSE si nRMSE pe train si test
return theta, metrici
\end{lstlisting}
\end{algorithm}

\section{Identificabilitate și limitări ale calibrării}
\label{sec:identificabilitate}

O eroare mică nu implică automat parametri epidemiologic corecți. În modelele
neliniare, mai multe combinații de parametri pot produce traiectorii asemănătoare
pe aceeași fereastră de timp. Această problemă se numește neidentificabilitate
practică și este accentuată atunci când datele observate sunt agregate.

Pentru SIDARTHE, neidentificabilitatea este mai importantă decât pentru SIR.
Dacă setul de date nu conține date despre testare, spitalizare, terapie intensivă și
vindecări, parametri precum $\epsilon$, $\zeta$, $\mu$, $\nu$, $\sigma$ sau
$\tau$ nu pot fi estimați robust doar din cazuri active. Într-o lucrare aplicată,
acest lucru trebuie declarat explicit. Rolul calibrării nu este să ascundă
incertitudinea, ci să arate ce ipoteze sunt necesare pentru a obține o potrivire
numerică.

\subsection{Vizualizarea identifiabilitatii practice pentru SIR}
\label{subsec:identif_landscape}

Pentru a face explicit acest fenomen, peisajul de pierdere
$\mathrm{MSE}(\beta, \gamma)$ a fost calculat pe o grilă $56 \times 36$ în
domeniul $\beta \in [0{,}05, 0{,}60]$, $\gamma \in [0{,}05, 0{,}40]$, pentru
fereastra de antrenare de mai sus
(Figura~\ref{fig:identifiability_landscape}). Reprezentarea în logaritm
relevă o vale alungită, aproape liniară în planul $(\beta, \gamma)$:
combinatii cu raport $\beta/\gamma$ aproximativ constant produc traiectorii
similare pe fereastra observată. În consecință, calibrarea liberă nu
distinge intre $(0{,}368, 0{,}335)$ și $(0{,}104, 0{,}0714)$ pe baza datelor
din valul 1 — informația lipsește în date, nu în algoritm.

\begin{figure}[t]
\centering
\begin{tikzpicture}
\begin{axis}[
    width=\columnwidth,
    height=7.5cm,
    xlabel={$\beta$ (zi$^{-1}$)},
    ylabel={$\gamma$ (zi$^{-1}$)},
    xmin=0.05, xmax=0.6,
    ymin=0.05, ymax=0.40,
    colormap/viridis,
    colorbar,
    colorbar style={
        title={$\log_{10}\mathrm{MSE}$},
        title style={font=\footnotesize},
        ticklabel style={font=\footnotesize},
    },
    point meta min=6.9,
    point meta max=10,
    tick label style={font=\footnotesize},
    label style={font=\small},
]
% peisajul de pierdere
\addplot[
    scatter, only marks, mark=square*, mark size=2.7pt,
    scatter src=explicit,
] table[col sep=comma, x=beta, y=gamma, meta=log_mse]
    {data/identifiability_grid_sir.csv};

% optimul liber
\addplot[mark=*, mark size=4pt, color=red, mark options={fill=white, line width=1pt}]
    coordinates {(0.3677, 0.3346)};
\node[font=\scriptsize, anchor=south, color=red, fill=white, fill opacity=0.85, text opacity=1, inner sep=1pt]
    at (axis cs:0.3677, 0.3346) {liber};

% optimul constrans (gamma=1/14)
\addplot[mark=*, mark size=4pt, color=red!70!black, mark options={fill=white, line width=1pt}]
    coordinates {(0.1039, 0.0714)};
\node[font=\scriptsize, anchor=south, color=red!70!black, fill=white, fill opacity=0.85, text opacity=1, inner sep=1pt]
    at (axis cs:0.1039, 0.0714) {constrans};

% linia gamma=1/14
\draw[dashed, red!50, thick] (axis cs:0.05, 0.0714) -- (axis cs:0.6, 0.0714);
\node[font=\scriptsize, anchor=west, color=red!70!black]
    at (axis cs:0.55, 0.0714) {$\gamma = 1/14$};

\end{axis}
\end{tikzpicture}
\caption{Peisajul de pierdere $\log_{10}\mathrm{MSE}(\beta, \gamma)$ pentru
modelul SIR pe fereastra de antrenare. Punctele rosii marcheaza solutia
libera $(0{,}368, 0{,}335)$ si solutia constransa $(0{,}104, 0{,}0714)$;
linia rosie punctata indica restrictia clinica $\gamma = 1/14$. Valea de
pierdere mica este alungita si aproape liniara, ceea ce inseamna ca multe
combinatii $(\beta, \gamma)$ produc o potrivire similara cu datele observate
din valul 1 --- semnatura tipica a identifiabilitatii practice slabe pentru
SIR cu o singura serie observata.}
\label{fig:identifiability_landscape}
\end{figure}


Această observație justifică restrictionarea $\gamma$ din literatura clinică
ca o practică recomandabilă, nu ca o limitare arbitrară. Mai mult, rezolvarea
unei astfel de identifiabilitati slabe pentru SIDARTHE necesită informatii
suplimentare în date — testare, spitalizare, terapie intensivă, decese — care
permit fixarea sau constrangerea unor parametri prin operatori de observație
multipli, nu doar serie singulară de cazuri active.

\section{Integrarea în sistemul software}
\label{sec:integrare_calibrare_software}

Arhitectura curentă este pregătită pentru calibrare deoarece separă trei
elemente: datele observate, parametrii modelului și rezultatele simulării.
Extinderea naturală constă într-un tip suplimentar de job, de exemplu
\texttt{calibration}, care rulează mai multe simulări candidate și salvează
parametrii rezultați într-o tabelă dedicată. Simularea finală ar putea fi apoi
creată automat folosind $\hat{\theta}$.

Această extensie trebuie totuși tratată diferit de o simulare obișnuită.
Calibrarea poate necesita zeci sau sute de integrări numerice, deci timpul de
execuție, costul CPU și consumul de memorie devin parte a rezultatului. Din acest
motiv, evaluarea numerică și evaluarea software trebuie raportate împreună:
modelul cu eroare mai mică poate fi mai costisitor, iar modelul mai rapid poate
fi insuficient pentru datele disponibile.

\testtodo{Pentru integrarea software a calibrării trebuie măsurate costurile
generate de rulări multiple ale solverului: timp total, timp CPU, memorie, număr
de job-uri, număr de simulări candidate și volum de rezultate scrise în baza de
date.}
